\subsection{Inference and the KV cache}
\label{sec:transformer-inference}

Inference uses the same parameters but does not know future tokens. It must generate one
token, append it to the sequence, and then generate the next. Efficient inference splits
this work into two phases.

\paragraph{Prefill.}
The entire prompt is known, so it is processed in parallel. At every layer, the engine
retains the key and value vectors produced for every prompt position.

\paragraph{Decode.}
After sampling a token, the model processes only that new position. Its query reads the
past keys and values, while its new key and value are appended to the cache. Past hidden
states do not need to be recomputed because the causal mask guaranteed that they could
not depend on the newly appended token.

Suppose a prompt has $P=4096$ tokens and the model generates $G=1024$ tokens. Without a
cache, producing each token reruns its complete available prefix:
\begin{equation}
  \sum_{g=0}^{G-1}(P+g)
  =GP+\frac{G(G-1)}{2}
  =4{,}718{,}080
\end{equation}
token positions through the model. With caching, prefill processes the 4096 prompt tokens
once. The first output is sampled from the final prefill logit, so only the first
$G-1$ generated tokens must be processed to obtain the remaining outputs. The total is
$P+G-1=5119$ token positions, about $921.7$ times fewer token-local evaluations in this
example.

This compute saving consumes memory. If a layer uses $H_{\mathrm{kv}}$ key/value heads,
head width $d_h$, current context length $T$, batch size $B$, and $s$ bytes per stored
number, its cache occupies
\begin{equation}
  2BH_{\mathrm{kv}}Td_hs\ \text{bytes}.
\end{equation}
The leading 2 is for keys and values. With one sequence, 32 key/value heads, head width
128, context length 5120, and two-byte values, the cache is 80 MiB per layer and 2.5 GiB
across 32 layers. Grouped-query and multi-query attention reduce the number of key/value
heads and therefore reduce this storage and the associated memory traffic
\citep{shazeer2019fastdecoding}.

KV caching removes repeated projections and MLP work for old tokens. It does not make a
decode step constant-cost: each new query must still read keys and values from a context
that grows with every generated token. Training does not normally use this decode-style
cache because all training tokens are already known and can be processed in parallel,
while backpropagation needs the intermediate activations rather than only past keys and
values.
